\documentclass[12pt,a4paper]{article}

% ------------------------------------------------
% Paquetes
% ------------------------------------------------
\usepackage[utf8]{inputenc}
\usepackage[T1]{fontenc}
\usepackage[spanish,es-nodecimaldot]{babel}
\usepackage{lmodern}
\usepackage{microtype}
\usepackage{geometry}
\usepackage{setspace}
\usepackage{parskip}
\usepackage{amsmath}
\usepackage{amssymb}
\usepackage{csquotes}
\usepackage{hyperref}

% ------------------------------------------------
% Configuración
% ------------------------------------------------
\geometry{
    top=2.7cm,
    bottom=2.7cm,
    left=3cm,
    right=3cm
}

\setstretch{1.15}

\hypersetup{
    colorlinks=true,
    linkcolor=black,
    urlcolor=black,
    citecolor=black
}

\title{\textbf{De la colonización de la alienación}}
\author{Vicente Domínguez Arca}
\date{}

% ------------------------------------------------
% Documento
% ------------------------------------------------
\begin{document}

\maketitle

No supone un atentado contra la opinión particular el que este autor manifieste que la inmensa mayoría de las personas han sentido, en algún momento, que se habían equivocado. Se trata de un pensamiento muy personal, en el que nos situamos retrospectivamente en el momento de tomar una decisión, permanecemos allí durante un tiempo prudencial y, en función del devenir de los acontecimientos, nos reafirmamos en nuestra particular decisión o la lamentamos.

Esta práctica es sorprendentemente común. Tanto que, en muchas situaciones, jugamos a ser otros, tomando otras decisiones, y empleamos las consecuencias de lo sucedido para señalar responsabilidades en los demás. Nada más lejos del mecanismo original, en el que legítimamente podemos situarnos desde nuestra propia referencia, pero difícilmente podemos argumentar que nos corresponde desplazar al prójimo para ocupar su posición en una ficción cuya finalidad es responsabilizarlo, normalmente, de un error que consideramos objetivo.

De algún modo, esta práctica supone una colonización del prójimo o, al menos, de una parte de sus dominios personales. Nos establecemos en una parte ficticia del prójimo que nosotros mismos diseñamos y «cultivamos» esa sección para obtener un determinado producto, que suele ser lo que consideramos una explicación. No se trata de comprender al otro desde su propia posición, sino de ocuparla imaginariamente desde la nuestra. Es una operación que podemos realizar con tal rapidez que apenas advertimos el desplazamiento. Otras formas de aproximación al prójimo, como la empatía, requieren sosiego, tiempo y dedicación intelectual.

Ese acto es, en el sentido que aquí nos ocupa, un fenómeno de alienación: provocamos que algo del otro pase a ser también nuestro, volviendo parcialmente ajeno para él aquello que inicialmente le pertenecía. La interpretación de su experiencia deja de ser exclusivamente suya, pues ahora existe también bajo una forma construida por nosotros. Y, dependiendo de las relaciones de coacción, sometimiento o autosometimiento que intervengan, esa apropiación puede llegar a comprometer, parcial o totalmente, el control que el prójimo conserva sobre su propia experiencia.

Aparece así una especie de relación simbiótica entre colonización y alienación.

Esta cadena de tráfico informativo puede encontrarse, y se encuentra, con mucha frecuencia. Son numerosos los ejemplos cotidianos en los que esta simple sucesión suele determinar el hilo conductor; tantos que no podemos hacer aquí una compilación de ellos, por prudencia y por intentar mantener un orden epistemológico. Porque la cuestión no es en cuántas situaciones aparece esta secuencia, sino qué lógica encierra dentro del problema de lo común o dentro del problema de lo social.

Son numerosos los esfuerzos actuales, en el ámbito de la teoría crítica, por identificar las bases del *burnout*, de la alienación y del modo en que esta se despliega en una sociedad occidental cuyos individuos viven endeudados. Encontramos las teorías de la aceleración de Hartmut Rosa, las reflexiones sobre el agotamiento de Byung-Chul Han o el análisis del sujeto endeudado de Maurizio Lazzarato, asociado a la figura del *homo debitor*.

Y surge un problema muy concreto: si esa colonización de la alienación conduce a disputas entre individuos, ¿qué provoca que no se extinga por improductividad esa tendencia?

Existen, a juicio de este autor, dos posibilidades que merecen ser examinadas. Una es que el ejercicio de esa colonización de un individuo sobre otro repercuta en algo positivo para el colonizador. La otra posibilidad es que exista un fondo común, un sustrato social que sostenga la continuidad de esta práctica mediante, claro está, su fomento.

La primera explicación puede parecer verosímil, pero, si uno interroga esa posibilidad en detalle, acaba por cerciorarse de que no es más que una desagradable ficción como explicación suficiente. Sencillamente, para que el colonizador acabe considerando positivo su ejercicio, debe existir un sustrato social en el que ese tipo de actitudes encuentre algún reconocimiento o recompensa. Y eso desemboca directamente en la segunda posibilidad.

Después de reflexionar lo que considero suficiente sobre esa posibilidad, he llegado a la conclusión de que es en una sociedad donde existen la desigualdad y la jerarquía donde esos ejercicios encuentran condiciones especialmente propicias para su fomento. La verticalidad no solo permite que una interpretación ajena adquiera mayor autoridad que la de quien ha vivido una experiencia, sino que puede favorecer que este último termine aceptándola como propia. Es una forma de coacción, o así podríamos entenderla, que no siempre necesita manifestarse mediante una imposición explícita. A su modo, alimenta la desigualdad y fortalece determinadas estructuras verticales.

No trato aquí de identificar directamente al neoliberalismo como el máximo responsable, ni siquiera como catalizador. Pero sí como oxígeno de la combustión indicada.

El propio eje del neoliberalismo concede, de modo tácito pero muy rentable, la posibilidad de asumir el papel de colonizador a todo aquel individuo que muestre capacidades para hacerlo. Tanto que podría considerarse una vía acertada para beneficiarse de las maquinarias reiterativas que constituyen el sistema neoliberal. Allí donde la posición relativa de los individuos adquiere importancia, la capacidad de interpretar las decisiones ajenas, atribuir responsabilidades y obtener reconocimiento por ello puede convertirse en un recurso especialmente provechoso.

La cuestión, después de todo esto, es simple.

Vivimos rodeados de innumerables ejemplos que pueden ser descritos mediante la dinámica que hemos presentado, pero no encuentro tan numerosos ejemplos de otras dinámicas posibles.

¿Hay alternativas a la colonización sistemática de la alienación?

¿Merece la pena afanarse en buscarlas e implementarlas?

¿Conduciría esa búsqueda a alguna vía de escape frente a presiones insostenibles?

\end{document}